% ===== 原文 pp.103--106 =====
% [tr-08 起点：本段接 tr-07 在 p.102 末句 "which are available provided the sum of" 之后；
%  若 tr-07 已把该句补完到 p.103 首行 "…is suitably quantized."，则本行起的"内部动量…"自然接续。]

% [原文 p.103]
共形块中的内部动量则对应矩阵模型积分中围道的选取 [289, 1081, 1082]。要离开参数空间里的这些量子化切片，就需要做解析延拓；而解析延拓只有在 $\fg=\fsu(2)$ 的情形下才完全受控，因为此时各种展开的系数已知是全部参数的有理函数。

% [Tutor Note: p.103 的 $b^2=\epsilon_1/\epsilon_2\ne-1$ 要这样读：未精细化（unrefined）点是 $\epsilon_1+\epsilon_2=0$，等价于 $b^2=-1$，即 $b=\pm\ii$（纯虚参数），对实 $b$ 根本达不到；这正是原文强调"离开量子化切片需要解析延拓、且只有 g = su(2) 情形完全受控"的原因。约定核对：原文在 p.8（式 (1.7) 之后）与 p.52（§5.3 Omega 背景处）都写明 $\epsilon_1=b/r$、$\epsilon_2=1/(rb)$，故 $\epsilon_1/\epsilon_2=b^2$；于是 Nekrasov--Shatashvili 极限 $\epsilon_2\to0$ 等价于 $b\to\infty$，这与附录 A 中 $\Gamma_b,S_b,\Upsilon_b$ 在 $b\to1/b$ 下不变（即 $\epsilon_1\leftrightarrow\epsilon_2$ 互换）相容------p.105 缩写表 "NS" 条写作"$b\to0$ 极限"，正是在 $b\to1/b$ 对称下对同一个 NS 极限的表述（等价地 $\epsilon_1\to0$）。]

4d 规范理论与矩阵模型之间的联系表现为：$Z_{\mathrm{inst}}$ 与矩阵模型配分函数相等 [1083, 1084]（不必绕道拓扑弦 $Z_{\mathrm{top}}$，而是直接相等），SW 曲线与矩阵模型的谱曲线相符、SW 微分与单点预解式（1-point resolvent）相符 [60, 370, 894, 1085, 1086]。这条联系以及与 2d CFT 的联系都可作如下推广：在 $\beta$ 形变矩阵模型的框架下推广到 $b^2=\epsilon_1/\epsilon_2\ne-1$ 的情形 [1083, 1086]，并推广到 5d [618, 1083, 1087]、渐近自由理论 [370, 373, 1088--1091]、$\fsu(N)$ 理论 [1087, 1092]、箭图规范理论 [1093] 与广义箭图 [1094]（此时曲面 $C$ 的亏格更高），还推广到 $\R^4$ 的轨形 [21]。

这些关系在 [373, 1095--1099] 中以各种极限做了检验，并在某些情形下得到证明 [67, 307]。由于 $Z_{\mathrm{inst}}$ 的 Nekrasov--Shatashvili 极限 $\epsilon_2\to0$ 会对 4d $\cN=2$ 理论 SW 解背后的可积模型做量子化，矩阵模型能提供关于可积模型的有用信息，例如见 [729, 919, 1100--1103]。矩阵模型也被研究过，用于野生穿刺与 AD 理论方面的应用，特别是在经典极限 $c\to\infty$（即规范理论一侧的 Nekrasov--Shatashvili 极限）下，见 [371, 372, 377, 379, 380, 384, 386--391, 1104--1108]。

另一个方向上，文献 [218, 219, 224, 225, 573, 996, 997, 1109--1120] 借助矩阵模型及其他技巧，研究了（适当归一化的）瞬子配分函数在 S 对偶下的模性质。相关的工作与综述还有很多 [301, 629, 639, 1121--1134]，此外还有一些博士论文 [1135--1137]。

\para{其他相关主题。}
我们在此罗列一些以各种方式与 AGT 对应相关联、彼此却颇为不同的主题。与综述其他部分相比，这里的文献列表既不那么完整，筛选得也没那么严格。

6d $\cN=(2,0)$ 理论、4d $\cN=2$ class S 理论，以及带延展算子（extended operators）的情形，其全息对偶分别在 [448, 557, 587, 763, 764, 766, 1052, 1138--1156] 中做了探讨。

超对称局域化适用于许多背景几何，感兴趣案例的一个样本见 [240, 503, 822, 848, 1157--1167] 以及综述 [1160, 1163]。$Z_{\mathrm{inst}}$ 的各种展开（以及超对称局域化得到的其他精确结果）的复活与 Borel 可和性在 [468, 514, 1168--1173] 中做了研究。这些工作对复活技巧在 QFT 中究竟有多大适用范围给出了一些洞见。

与 4d $\cN=2$ 理论的其他性质之间还有众多相互关联。

\begin{itemize}[leftmargin=1.4em,itemsep=3pt,topsep=3pt]
  \item 拓扑--反拓扑融合（$tt^*$ 方程）[999, 1174--1176]。

% [Tutor Note: p.103 底部与 p.104 顶部的 bullet 同属一个 itemize（p.103 末项为"手征代数"那条），跨页断开由排版自动处理。$tt^*$ 方程指 Cecotti--Vafa 的拓扑--反拓扑融合方程（$t$ 与 $\bar t$ 平面上的 tt* 方程），它等价于 4d $\cN=2$ 理论特殊几何（special Kähler 结构）的那组关系；术语首现：topological anti-topological fusion = 拓扑--反拓扑融合。]

  \item 作为 4d $\cN=2$ 理论受保护子部门（protected subsector）出现的手征代数 [172, 423, 580, 681, 683, 792, 795, 949, 1067, 1177--1217]（另见 [293, 1218--1222] 及其中的文献）。

% [原文 p.104]
  \item 规范/Yang--Baxter 方程（YBE）[674, 713, 1223--1229]，综述见 [1230]，另见 [1231]。
  \item class S 理论是由基本构件通过规范化（gauging）拼装而成的，这一构造方式提示我们引入理论空间（theory space）的概念 [1232, 1233]。
  \item 共形自举（conformal bootstrap）：除了本综述讨论的构造之外，另一种寻找 QFT------特别是幺正 CFT------的有趣方法，是始于 [1234]、并被应用于 4d $\cN=2$ 理论的共形自举纲领 [1235]。特别地，某些 AD 理论恰好位于共形自举所允许的参数空间区域的角上。
\end{itemize}

AGT 对应提升了人们对 2d CFT 中若干老问题的兴趣。

\begin{itemize}[leftmargin=1.4em,itemsep=3pt,topsep=3pt]
  \item 计算共形块、关联函数与融合矩阵：或通过递推关系 [227, 294, 371, 619, 1118, 1236--1244]，或利用大 $c$ 极限下的全息 [298, 479, 891, 1245--1262]，或用 Chern--Simons 理论 [1263]。
  \item 研究 Toda CFT 的变体、仲费米子 Liouville CFT 等 [1264--1266]。
  \item 与分数量子霍尔效应的一些（有争议的）联系 [1141, 1267--1272]。
  \item 等单值问题（isomonodromy problems）：现已知道，关于内部动量做 Fourier 变换后的共形块（从而是 Nekrasov 配分函数）给出 Fuchs 型联络的等单值问题中出现的 Painlevé 方程的解 [233, 394, 398, 401, 591, 910, 925, 946, 989, 1002, 1106, 1120, 1130, 1131, 1169, 1173, 1246, 1259, 1273--1305]；同样，$q$-形变 Virasoro 代数与 $q$W 代数的共形块给出 $q$-Painlevé 方程的解 [636--638]。
\end{itemize}

顺带一提，Liouville CFT 终于从其路径积分出发获得了数学上的严格定义：见 [1306, 1307] 及其中的文献。其他数学文献还包括对（$U_q(\mathfrak{sl}_2)$ 之模双（modular double））的 6j 符号的研究 [88]，以及与几何 Langlands 对应或其形变之间的联系 [517, 644, 1063, 1064, 1308, 1309]。

\para{最后的思考。}
在各种几何设置中做维数约化来构造新理论，已被证明极富成效。它产生了大量新的量子场论，这些理论可以作为构件，去得到更多的发现。用这种方式发现的众多对偶，还可以通过考虑各种形态的延展算子而进一步丰富。我希望读者们会参与这段令人兴奋的旅程，一起去描绘量子场论的版图！

% [Tutor Note: p.104 原文 "The construction of new theories by dimensionally reduction" 中 "dimensionally reduction" 应是 "dimensional reduction"（维数约化）之笔误：判断依据是全篇一律用 "dimensional reduction / compactification"（如摘要与 §9 的用法），且 "dimensionally" 是副词、不能修饰名词 reduction。译文按正确含义处理。]
